% ------------------------------------------------------------
% BORRADOR: Dataset y metodología adaptado a criptomonedas
% Reemplazo/complemento para Actividad 6
% ------------------------------------------------------------

% ------------------------------------------------------------
\subsection{Dataset y metodología}
% ------------------------------------------------------------

Se analiza una serie de alta frecuencia del mercado de criptoactivos,
específicamente los retornos logarítmicos minuto a minuto del par
BTC/USDT (Bitcoin vs.~Tether) obtenidos mediante la API pública de
Binance, el exchange de criptomonedas con mayor volumen operativo
del mundo\footnote{Binance API v3, endpoint \texttt{/api/v3/klines}.
Documentación oficial: \url{https://binance-docs.github.io/apidocs/spot/en/}}.

La elección de Bitcoin como objeto de estudio responde a cuatro propiedades
estructurales que lo hacen ideal para el análisis de fractalidad:

\begin{enumerate}[leftmargin=2em]
    \item \textbf{Dependencia de rango largo documentada.}
          Múltiples estudios empíricos reportan que los retornos de alta
          frecuencia de Bitcoin exhiben memoria persistente con
          $H \approx 0{,}60$--$0{,}70$
          \cite{bariviera2017inefficiency,dionisio2022long},
          situándose en un régimen intermedio entre procesos sin memoria
          ($H=0{,}5$) y sistemas biológicos de alta LRD ($H\approx0{,}96$).

    \item \textbf{Alta frecuencia de muestreo.}
          La API proporciona velas (candles) con granularidad de 1 minuto,
          lo que permite construir series de miles de observaciones en
          ventanas temporales de pocos días, superando la resolución del
          dataset NASA HTTP, cuya unidad nativa es la solicitud individual
          con tiempos entre llegadas altamente irregulares.

    \item \textbf{Mercado continuo 24/7.}
          A diferencia de los mercados bursátiles tradicionales, Bitcoin
          no tiene cierres diarios ni fines de semana, eliminando los
          \textit{overnight gaps} que introducen discontinuidades artificiales
          en series temporales financieras clásicas y facilitando la
          aplicación de estimadores fractales que asumen regularidad
          muestral.

    \item \textbf{Alta liquidez y profundidad de mercado.}
          El par BTC/USDT registra volúmenes superiores a $10^9$ USD diarios,
          garantizando que los precios reflejen descubrimiento dinámico
          genuino y no artefactos de baja liquidez o manipulación de
          spreads.
\end{enumerate}

La serie objeto de estudio son los retornos logarítmicos de los precios de
cierre:
\begin{equation}
    r_t = \ln\!\left(\frac{P_t}{P_{t-1}}\right),
    \qquad t = 1, \dots, n,
\end{equation}
donde $P_t$ es el precio de cierre de la vela en el minuto $t$.
Esta transformación estabiliza la varianza, centra la serie en torno a
cero y es el estándar en econometría financiera para el análisis de
memoria persistente.

Se descargaron $n = 5\,000$ velas de intervalo 1m mediante la API pública
de Binance, abarcando el periodo del 24 de mayo de 2026 (17:33 UTC) al
28 de mayo de 2026 (04:52 UTC), equivalente a aproximadamente 3 días y
11 horas de operación continua. La serie resultante contiene $n-1 = 4\,999$
retornos logarítmicos con las siguientes estadísticas descriptivas:
media $\mu = -9{,}32 \times 10^{-6}$, mediana $0{,}0000$,
desviación estándar $\sigma = 4{,}5394 \times 10^{-4}$,
mínimo $-5{,}3145 \times 10^{-3}$, máximo $3{,}1934 \times 10^{-3}$,
curtosis $9{,}83$ y asimetría $-0{,}35$. La curtosis leptocúrtica y la
col pesada negativa indican la presencia de eventos extremos (\textit{flash
crashes} y \textit{pumps}) que constituyen candidatos naturales para
heterogeneidad de escalamiento multiescala.

\begin{table}[H]
    \centering
    \caption{Estadísticas descriptivas de los retornos logarítmicos
             BTC/USDT (1m).}
    \label{tab:btc-stats}
    \begin{tabular}{@{}lr@{}}
        \toprule
        \textbf{Estadística} & \textbf{Valor} \\
        \midrule
        Observaciones ($n$) & 4\,999 \\
        Media ($\mu$) & $-9{,}32 \times 10^{-6}$ \\
        Mediana & $0{,}0000$ \\
        Desv.~est.~($\sigma$) & $4{,}5394 \times 10^{-4}$ \\
        Mínimo & $-5{,}3145 \times 10^{-3}$ \\
        Máximo & $3{,}1934 \times 10^{-3}$ \\
        Asimetría & $-0{,}3470$ \\
        Curtosis & $9{,}8258$ \\
        Precio inicial & $76\,627{,}52$ USDT \\
        Precio final & $73\,139{,}71$ USDT \\
        \bottomrule
    \end{tabular}
\end{table}

Los estimadores VT y LD, así como el clasificador mono/multifractal de la
Actividad~4, se aplican sin modificación metodológica, garantizando
comparabilidad con los dominios previos. Esta continuidad es esencial:
permite situar la serie Bitcoin en el mismo espacio de parámetros que las
secuencias sintéticas, las señales cardíacas y el tráfico web NASA.

Desde la perspectiva de la taxonomía transversal desarrollada en este
trabajo, Bitcoin ocupa una posición estratégicamente intermedia. El tráfico
NASA HTTP presenta memoria persistente modesta ($H \approx 0{,}54$) pero
heterogeneidad espacial pronunciada (multifractal), mientras que la
fisiología cardíaca exhibe LRD intensa ($H \approx 0{,}96$) con marcada
multifractalidad. Los retornos de Bitcoin, con $H$ esperado en el rango
$0{,}60$--$0{,}70$, llenan el vacío conceptual entre ambos extremos:
un sistema tecnológico--financiero con memoria persistente intermedia que
no llega a la complejidad adaptativa del sistema nervioso autónomo, pero
que supera ampliamente la estructura de llegadas de Poisson asumida en
modelos clásicos de tráfico web. Esta ubicación intermedia fortalece la
hipótesis transversal: la fractalidad no es epifenómeno de un dominio
particular, sino propiedad emergente de sistemas complejos con
retroalimentación, multiescala y heterogeneidad estructural.
